\section{Human-Curated SFT Results: el evaluator cambia la conclusión}

Ahora volvemos al experimento con los datos human-written de Surge. La clase va cambiando, uno tras otro, el dataset source, la prompt/response length y el dataset size, y observa al mismo tiempo las Open LLM metrics y MT-Bench. El resultado más importante no es cuál barra queda más alta, sino que los evaluadores con frecuencia producen ordenamientos distintos.

\subsection{El conflicto entre las Open LLM Metrics y MT-Bench}

La sección anterior ya explicó que cada evaluator responde una pregunta distinta; esta sección usa el mismo grupo de checkpoints de LLaMA-2-13B para comprobar si esa diferencia realmente cambia el dataset ranking. Al leer las dos gráficas de barras, primero se fijan el tamaño del modelo y la etapa de entrenamiento, y después se compara si el apparent winner cambia cuando solo se sustituye el evaluator. El Open LLM Leaderboard se inclina por las task metrics estándar, mientras que MT-Bench se inclina por la calificación que el judge asigna a respuestas de varios turnos; como sus ponderaciones son distintas, es natural que el ranking de los modelos se invierta.

\lecturefigure{slide-49.jpg}{Resultados de distintos SFT datasets en el Open LLM Leaderboard}{Diapositiva oficial p.49}

\begin{knowledgebox}{Lectura de la figura: primero cada benchmark, después el aggregate}
La gráfica de barras se compone de varios task scores, y un aggregate alto puede deberse a un solo punto fuerte. Primero se comparan los subcomponentes, como TruthfulQA y reasoning/knowledge, y después se revisa la task mix del conjunto de datos; la figura muestra correlación y no demuestra por sí sola que una fuente de datos cause una mejora de capacidad.
\end{knowledgebox}

\lecturefigure{slide-50.jpg}{MT-Bench scores del mismo grupo de SFT models}{Diapositiva oficial p.50}

\begin{warningbox}{La metric reversal es una measurement signal}
Si Open LLM y MT-Bench eligen ganadores distintos, no conviene escoger el resultado que más guste. Lo correcto es localizar la diferencia: ¿MT-Bench prefiere respuestas más largas? ¿Las metrics estándar ignoran el estilo conversacional? ¿Algún conjunto de datos se ajusta mejor al judge? El conflicto mismo pone al descubierto la evaluation function.
\end{warningbox}

\teachervoice{La ponente señala que las conclusiones de las metrics automáticas y de MT-Bench a veces son «opposite». No declara que una de las tablas de clasificación sea la verdad, sino que lleva la contradicción a los análisis posteriores de length y del sesgo de GPT-4.}

\subsection{La Response Length como Confounder}

En la subsección anterior se observó la inversión del ranking; esta subsección rastrea el confounder que con más facilidad se pasa por alto: los distintos conjuntos de entrenamiento moldean modelos con distinta longitud y formato de respuesta. Al leer la tabla de longitudes, primero hay que tratarla como una style statistic y no como una calificación de calidad, y después observar si el judge podría premiar esos rasgos superficiales. La average completion length varía mucho entre datasets; si el judge prefiere textos con explicaciones completas, en forma de lista o más largos, la length puede influir a la vez en el style del entrenamiento y en la calificación de la evaluación, lo que produce la ilusión de algo «más helpful».

\lecturefigure{slide-51.jpg}{Response length promedio de Surge, LIMA y OpenAssistant}{Diapositiva oficial p.51}

\begin{knowledgebox}{Lectura de la figura: 211, 482 y 722 son un style proxy, no una calificación de calidad}
Surge es más breve, LIMA queda en medio y OpenAssistant es más largo. Las cifras solo describen la longitud promedio; para juzgar la calidad también hay que controlar la task, la exactitud factual y la redundancia. Tomar automáticamente una respuesta larga como más completa premia la verbosity y no la usefulness.
\end{knowledgebox}

\subsection{Prompt Length y Performance}

La clase observa además varias metrics según la average prompt length. Un prompt largo puede contener más contexto, pero también puede corresponder a una tarea más difícil; por eso la dirección de la curva debe interpretarse junto con la task composition.

\lecturefigure{slide-52.jpg}{Relación entre la performance y la average prompt length}{Diapositiva oficial p.52}

\begin{knowledgebox}{Lectura de la figura: en curvas de varios paneles, primero la consistencia}
Primero se comprueba si las pendientes de cada benchmark van en la misma dirección y después se buscan el turning point y las líneas atípicas. Si solo sube una metric, la length podría estar mejorando la behavior que esa metric prefiere, no la capacidad en general. Las curvas tampoco controlan la task difficulty.
\end{knowledgebox}

\lecturefigure{slide-53.jpg}{MT-Bench scores con distintos prompt-length settings}{Diapositiva oficial p.53}

\begin{warningbox}{La reacción del LLM judge ante la length no es estable}
En clase se encontró que MT-Bench y parte de las metrics automáticas vuelven a no coincidir, y que GPT-4 no siempre prefiere sin más el modelo asociado a prompts más largos. El judge bias es multidimensional: la verbosity de la respuesta, el formato, el estilo de entrenamiento y la task actúan en conjunto.
\end{warningbox}

\subsection{Dataset Size Ablation y Diminishing Returns}

El último grupo de experimentos con human data extrae subconjuntos de distinto tamaño de Surge-Instruct. La ablation se acerca más a una relación causal que la comparación entre conjuntos de datos, porque la provenance y la task taxonomy son más homogéneas, aunque todavía hay que considerar la varianza del muestreo, los training tokens y el evaluator.

\lecturefigure{slide-54.jpg}{Indicadores automáticos de la dataset size ablation de Surge-Instruct}{Diapositiva oficial p.54}

\begin{knowledgebox}{Lectura de la figura: la pendiente marginal y la incertidumbre}
A medida que aumentan las muestras, varias curvas tienden a una meseta, lo que respalda los diminishing returns. La figura no muestra todas las semillas aleatorias ni los intervalos de confianza, así que no conviene sobreinterpretar diferencias pequeñas; lo más confiable es la tendencia de que «la mayor parte de la ganancia aparece pronto y seguir agregando datos similares rinde cada vez menos».
\end{knowledgebox}

\lecturefigure{slide-55.jpg}{Resultados de MT-Bench para la misma dataset-size ablation}{Diapositiva oficial p.55}

\begin{knowledgebox}{Lectura de la figura: el evaluator también reordena la misma ablation}
Los cambios en MT-Bench no tienen por qué ir en la misma dirección que los indicadores automáticos. Si agregar datos mejora las tasks estándar pero no el judge score, puede tratarse de una separación entre la task coverage y el conversational style; a la inversa, también puede tratarse de una judge style reward. La conclusión del experimento debe expresarse como un metric vector, no como una calificación total.
\end{knowledgebox}

\teachervoice{La conclusión de Rajani es que, entre las metrics automáticas, TruthfulQA es la que más diferencia a los modelos, mientras que MT-Bench tiene una correlación limitada con esas metrics. Ella interpreta este fenómeno como un benchmark gap, en lugar de declarar sin más que MT-Bench o los benchmarks tradicionales fracasan.}

\subsection{Resumen de la sección}

En el Human-curated SFT, la fuente, la longitud y el tamaño de los datos influyen en el modelo, pero el evaluator cambia de manera notable el apparent winner. Un experimento de alta calidad debe reportar a la vez las task metrics, los judge scores, las response statistics y la human evidence, y tratar las discrepancias como resultados que hay que explicar.
